% Historical source (removed by author request): references/_OLD_Übungsaufgaben/Übungsblatt 7 - Komplexe Zahlen/Blatt7.tex:87
\begin{ExerciseBlock}
\ExerciseHeading{5}{Geometrische Interpretation}{complex_numbers_geometric_interpretation}
Was passiert geometrisch mit einer komplexen Zahl $z$ durch Anwendung der folgenden Funktion?

\begin{enumerate}[label=\alph*)]
	\item $f(z) = z+1$
	\item $f(z) = z+i$	
	\item $f(z) = z \cdot i$
	\item $f(z) = z/i$
	\item $f(z) = -z$
	\item $f(z) = z^2$
	\item $f(z) = 1/z$
	\item $f(z) = \frac{z+z^*}{2}$
\end{enumerate}
Hinweis: Verwenden Sie für Ihre Überlegung eine sinnvolle Darstellung der komplexen Zahl.

\end{ExerciseBlock}
\UseExerciseSolution{complex_numbers_geometric_interpretation}
